\section{Conclusion}
\label{sec:conclusion}

We asked whether a hybrid Bayesian--symbolic framework can distinguish genuine
robustness from a weak intervention and from a design flaw as the cause of a null
result in an alignment experiment, more accurately than either modality alone. The
answer is yes. Fusing Bayesian evidence and power quantification with symbolic
reasoning about experimental structure classifies the cause of a null at
\macrofone $0.991$ on a held-out benchmark, far above Bayesian-only ($0.489$),
symbolic-only ($0.871$), naive NHST ($0.244$), equivalence-test-only ($0.326$), and
a frontier LLM given identical information ($0.888$). The reason is mechanistic:
fusion is the only approach that observes both the statistical-adequacy signals and
the structural/identifiability signals, and it delivers its verdicts as transparent,
actionable reasoning chains.

\para{Key takeaway.} An evidence layer and a structure layer are each necessary but
individually blind---to design flaws and to power respectively---and only their
fusion can read a null correctly.

\para{Future work.} We see five concrete next steps. (1) Validate on semi-synthetic
scenarios anchored on real \harmbench{} attack-success structure, and ultimately on
a small human-expert-labelled corpus of published alignment nulls. (2) Replace the
fixed thresholds with per-study power analyses and a measurement-validity audit.
(3) Post-hoc calibrate the confidences and emit partial-cause (multi-label)
decompositions for nulls with compound causes. (4) Convene a human-expert panel to
measure inter-annotator $\kappa$. (5) Stress-test threshold sensitivity as
thoroughly as we tested prior sensitivity here.
